% GENERATED FILE. Edit canonical lessons, then run npm run generate:books.
% canonical-source-hash: fnv1a64:d5ac521cdad7c1ea
% canonical-lessons: AR-C68-burtuqal, AR-C68-laymun, AR-C68-bayd, AR-C68-basal, AR-C68-futur

\chapter{Oranges, Lemons, Eggs, Onions, Breakfast}
\label{ch:ar-oranges-lemons-eggs-onions-breakfast}
\hlchaptermodality{68}

\begin{chapteropening}
\textbf{By the end of this chapter:} \emph{I can say oranges, lemons, eggs, onions and breakfast.}

Complete the last lesson of chapter 68: I can say oranges, lemons, eggs, onions and breakfast.
\end{chapteropening}

\section[\texorpdfstring{\ar{برتقال} (burtuqāl)}{برتقال (burtuqāl)}]{\ar{برتقال} (burtuqāl) — oranges}

\label{lesson:AR-C68-burtuqal}



\begin{quote}
Before the new one: say the Arabic for honey, then the Arabic for grapes.
\end{quote}

\subsection*{You'll want to know: \ar{برتقال}}
\textbf{\ar{برتقال}} — \emph{burtuqāl} — \textquotedblleft{}oranges\textquotedblright{}.

\begin{cousinweb}
\textbf{\ar{برتقال}} (\emph{burtuqāl}) is named after Portugal, whose traders are usually said to have brought the sweet orange from the east; EtymDB takes the Arabic from Greek \emph{portokali}, and Turkish has \emph{portakal}. It has no Arabic root.
\end{cousinweb}

\begin{grammarlens}[title={What you've built}]
One more word for this chapter.
\end{grammarlens}

\subsection*{Your turn}
\begin{itemize}
\raggedright
  \item \emph{Say it:} \emph{burtuqāl}
  \item \emph{Say it:} \emph{burtuqāl}, once more
  \item \emph{Say it:} say \emph{burtuqāl}
  \item \emph{Recall:} say the Arabic for honey, then the Arabic for grapes, then say \emph{burtuqāl} again
\end{itemize}

\begin{tcolorbox}[breakable,colback=teal!4,colframe=teal!35!black,title={Before you move on}]
What does \ar{برتقال} mean? (\textquotedblleft{}Oranges\textquotedblright{}.) Say it once more.
\end{tcolorbox}
\section[\texorpdfstring{\ar{ليمون} (laymūn)}{ليمون (laymūn)}]{\ar{ليمون} (laymūn) — lemons}

\label{lesson:AR-C68-laymun}



\begin{quote}
Before the new one: say the Arabic for grapes, then the Arabic for oranges.
\end{quote}

\subsection*{You'll want to know: \ar{ليمون}}
\textbf{\ar{ليمون}} — \emph{laymūn} — \textquotedblleft{}lemons\textquotedblright{}.

\begin{cousinweb}
\textbf{\ar{ليمون}} (\emph{laymūn}) is a loan from Persian \emph{līmūn}, \emph{līmū}, \textquotedblleft{}lemon\textquotedblright{} (EtymDB), a word usually traced further east to an Indic source. It has no Arabic root.
\end{cousinweb}

\begin{grammarlens}[title={What you've built}]
One more word for this chapter.
\end{grammarlens}

\subsection*{Your turn}
\begin{itemize}
\raggedright
  \item \emph{Say it:} \emph{laymūn}
  \item \emph{Say it:} \emph{laymūn}, once more
  \item \emph{Say it:} say \emph{laymūn}
  \item \emph{Recall:} say the Arabic for grapes, then the Arabic for oranges, then say \emph{laymūn} again
\end{itemize}

\begin{tcolorbox}[breakable,colback=teal!4,colframe=teal!35!black,title={Before you move on}]
What does \ar{ليمون} mean? (\textquotedblleft{}Lemons\textquotedblright{}.) Say it once more.
\end{tcolorbox}
\section[\texorpdfstring{\ar{بيض} (bayḍ)}{بيض (bayḍ)}]{\ar{بيض} (bayḍ) — eggs}

\label{lesson:AR-C68-bayd}



\begin{quote}
Before the new one: say the Arabic for oranges, then the Arabic for lemons.
\end{quote}

\subsection*{You'll want to know: \ar{بيض}}
\textbf{\ar{بيض}} — \emph{bayḍ} — \textquotedblleft{}eggs\textquotedblright{}.

\begin{cousinweb}
\textbf{\ar{بيض}} (\emph{bayḍ}) is built on the native Arabic root \emph{b-y-ḍ}, \textquotedblleft{}to be white, to whiten\textquotedblright{} (Lane). \emph{Bayḍ}, \textquotedblleft{}eggs\textquotedblright{}, is named on the root of whiteness, which also gives \emph{abyaḍ}, \textquotedblleft{}white\textquotedblright{}.
\end{cousinweb}

\begin{grammarlens}[title={What you've built}]
One more word for this chapter.
\end{grammarlens}

\subsection*{Your turn}
\begin{itemize}
\raggedright
  \item \emph{Say it:} \emph{bayḍ}
  \item \emph{Say it:} \emph{bayḍ}, once more
  \item \emph{Say it:} say \emph{bayḍ}
  \item \emph{Recall:} say the Arabic for oranges, then the Arabic for lemons, then say \emph{bayḍ} again
\end{itemize}

\begin{tcolorbox}[breakable,colback=teal!4,colframe=teal!35!black,title={Before you move on}]
What does \ar{بيض} mean? (\textquotedblleft{}Eggs\textquotedblright{}.) Say it once more.
\end{tcolorbox}
\section[\texorpdfstring{\ar{بصل} (baṣal)}{بصل (baṣal)}]{\ar{بصل} (baṣal) — onions}

\label{lesson:AR-C68-basal}



\begin{quote}
Before the new one: say the Arabic for lemons, then the Arabic for eggs.
\end{quote}

\subsection*{You'll want to know: \ar{بصل}}
\textbf{\ar{بصل}} — \emph{baṣal} — \textquotedblleft{}onions\textquotedblright{}.

\begin{cousinweb}
\textbf{\ar{بصل}} (\emph{baṣal}) is built on the native Arabic root \emph{b-ṣ-l}, \textquotedblleft{}onions\textquotedblright{} (Lane). Hebrew \emph{betsel}, \textquotedblleft{}onion\textquotedblright{}, is its cousin.
\end{cousinweb}

\begin{grammarlens}[title={What you've built}]
One more word for this chapter.
\end{grammarlens}

\subsection*{Your turn}
\begin{itemize}
\raggedright
  \item \emph{Say it:} \emph{baṣal}
  \item \emph{Say it:} \emph{baṣal}, once more
  \item \emph{Say it:} say \emph{baṣal}
  \item \emph{Recall:} say the Arabic for lemons, then the Arabic for eggs, then say \emph{baṣal} again
\end{itemize}

\begin{tcolorbox}[breakable,colback=teal!4,colframe=teal!35!black,title={Before you move on}]
What does \ar{بصل} mean? (\textquotedblleft{}Onions\textquotedblright{}.) Say it once more.
\end{tcolorbox}
\section[\texorpdfstring{\ar{فطور} (fuṭūr)}{فطور (fuṭūr)}]{\ar{فطور} (fuṭūr) — breakfast}

\label{lesson:AR-C68-futur}



\begin{quote}
Before the new one: say the Arabic for eggs, then the Arabic for onions.
\end{quote}

\subsection*{You'll want to know: \ar{فطور}}
\textbf{\ar{فطور}} — \emph{fuṭūr} — \textquotedblleft{}breakfast\textquotedblright{}.

\begin{cousinweb}
\textbf{\ar{فطور}} (\emph{fuṭūr}) is built on the native Arabic root \emph{f-ṭ-r}, \textquotedblleft{}to cleave; to break a fast\textquotedblright{} (Lane). Lane gives \emph{fuṭūr} as the food or drink with which a fast is broken.
\end{cousinweb}

\begin{grammarlens}[title={What you've built}]
One more word for this chapter.
\end{grammarlens}

\subsection*{Your turn}
\begin{itemize}
\raggedright
  \item \emph{Say it:} \emph{fuṭūr}
  \item \emph{Say it:} \emph{fuṭūr}, once more
  \item \emph{Say it:} say \emph{fuṭūr}
  \item \emph{Recall:} say the Arabic for eggs, then the Arabic for onions, then say \emph{fuṭūr} again
\end{itemize}

\begin{tcolorbox}[breakable,colback=teal!4,colframe=teal!35!black,title={Before you move on}]
What does \ar{فطور} mean? (\textquotedblleft{}Breakfast\textquotedblright{}.) Say it once more.
\end{tcolorbox}
